\documentclass[10pt,a4paper]{amsart}
\usepackage[margin=19mm]{geometry}
\usepackage{amsmath,amssymb,mathtools}
\usepackage[scr=rsfs,cal=euler]{mathalfa}
\usepackage{xfrac}
\usepackage[unicode,hidelinks,bookmarks=false]{hyperref}
\newcommand{\ie}{i.e.\ }
\newcommand{\cf}{cf.\ }
\newcommand{\resp}{respectively\ }
\newcommand{\Subsection}{\subsection}
\providecommand{\nobreakdash}{\nobreak\hbox{-}}
\setlength{\emergencystretch}{2em}
\multlinegap0pt
\usepackage{amssymb}
\usepackage{enumerate}
\newtheorem{theorem}{Theorem}
\newtheorem{lemma}{Lemma}
\def\XXint#1#2#3{{\setbox0=\hbox{$#1{#2#3}{\int}$}
     \vcenter{\hbox{$#2#3$}}\kern-.5\wd0}}
\def\Xint#1{\mathchoice
   {\XXint\displaystyle\textstyle{#1}}%
   {\XXint\textstyle\scriptstyle{#1}}%
   {\XXint\scriptstyle\scriptscriptstyle{#1}}%
   {\XXint\scriptscriptstyle\scriptscriptstyle{#1}}%
   \!\int}
\def\aver#1{\Xint-_{#1}}
\newcommand{\meanint}{{\int{\mkern-19mu}-}}
\title{The Green Function for Elliptic Systems in the Upper-Half Space}
\author{Martin Dindo\v{s}, Dorina Mitrea, Irina Mitrea, Marius Mitrea}
\date{Source: arXiv:2603.11251v1 (CC BY 4.0)}
\begin{document}
\maketitle

\noindent\emph{Source and licence.} The Green Function for Elliptic Systems in the Upper-Half Space. Version arXiv:2603.11251v1 (CC BY 4.0), distributed by arXiv under the Creative Commons Attribution 4.0 International licence, http://creativecommons.org/licenses/by/4.0/, as recorded in the arXiv licence metadata for this exact version and shown on the abstract page https://arxiv.org/abs/2603.11251v1. Full licence notice: \texttt{licenses/arxiv-2603.11251-CC-BY-4.0.txt}.\par\smallskip

\noindent\textit{Editorial scope.} Counted unit: the article's lemma Lgav-TeD (source0.tex lines 1539--1607). The article's complete original proof of it is delivered with it (source0.tex lines 1609--2063, one \texttt{proof} environment of 455 lines). No mathematics is written, repaired or re-derived: the statement and the proof are the article's own lines, in the article's own notation, and no retained line is substituted or re-flowed.  Retained uncounted as the source-local context the proof needs: lines 207--209; lines 228--231; lines 960--963; lines 1379--1391; lines 1507--1534.  Nothing was omitted from any retained span, so the package carries no omission record.  No retained line contains a citation, so the wrapper carries no bibliography and no bibliography entry.  No retained line contains a figure environment, an image-inclusion command or drawing code, so no image is delivered and the package declares no asset dependency.  Editorial formatting only, in addition: an AMS \texttt{amsart} preamble that restates the article's own declarations and macros used by the retained lines, namely the environments \texttt{theorem}, \texttt{lemma} on separate counters, together with the standard packages those lines need.  The retained environments print in this document as Theorem 1, Lemma 1 in order of appearance, whereas the article prints the same items with the numbering of its own full document.  The complete native source of the article as distributed in the arXiv e-print for this exact version is bundled unchanged as \texttt{sources/196203/source0.tex}.
\begin{equation}\label{L-def}
L:=\Big(a^{\alpha\beta}_{rs}\partial_r\partial_s\Big)_{1\leq\alpha,\beta\leq M},
\end{equation}

\begin{equation}\label{Def-ELLa.INTR}
{\rm det}\Big[\big(a^{\alpha\beta}_{rs}\xi_r\xi_s\big)_{1\leq\alpha,\beta\leq M}\Big]\not=0
\,\,\text{ for every }\,\,\xi=(\xi_r)_{1\leq r\leq n}\in{\mathbb{R}}^n\setminus\{0\}.
\end{equation}

\begin{equation}\label{NT-Fct.23}
\big({\mathcal{N}}^E_\kappa u\big)(x'):=\|u\|_{L^\infty(\Gamma_\kappa(x')\cap E)},\qquad
\forall\,x'\in\partial{\mathbb{R}}^{n}_{+}\equiv{\mathbb{R}}^{n-1}.
\end{equation}

\begin{theorem}\label{ker-sbav}
Consider a homogeneous, constant coefficient, second-order, system $L$ satisfying the weak 
ellipticity condition ${\rm det}\,[L(\xi)]\not=0$ for each $\xi\in{\mathbb{R}}^n\setminus\{0\}$.
Then for each null-solution $u$ of $L$ in a ball $B(x,R)$ {\rm (}where $x\in{\mathbb{R}}^n$ and $R>0${\rm )}, 
$p\in(0,\infty)$, $\lambda\in(0,1)$, $\ell\in{\mathbb{N}}_0$, and $r\in(0,R)$, one has
%
\begin{equation}\label{detraz}
\sup_{z\in B(x,\lambda r)}|(\nabla^\ell u)(z)|
\leq\frac{C}{r^\ell}\left(\aver{B(x,r)}|u|^p\,d{\mathscr{L}}^n\right)^{1/p},
\end{equation}
%
where $C=C(L,p,\ell,\lambda,n)>0$ is a finite constant.
\end{theorem}

\begin{theorem}\label{theor:div-thm}
Assume that $K\subset{\mathbb{R}}^n_{+}$ is a compact set and that
$\vec{F}\in\big[L^1_{\rm loc}({\mathbb{R}}^n_{+})\big]^n$ is a vector
field satisfying the following conditions {\rm (}for some aperture parameter $\kappa>0${\rm )}:
%
\begin{enumerate}\itemsep=0.2cm
%
\item [(a)] ${\rm div}\,\vec{F}\in\mathcal{E}'_K({\mathbb{R}}^n_{+})+L^1({\mathbb{R}}^n_{+})$,
where the divergence is taken in the sense of distributions;
%
\item[(b)] the nontangential maximal function ${\mathcal{N}}_\kappa^{\,{\mathbb{R}}^n_{+}\setminus K}\vec{F}$ 
belongs to $L^1({\mathbb{R}}^{n-1})$;
%
\item [(c)] the nontangential boundary trace 
$\vec{F}\big|_{\partial{\mathbb{R}}^n_{+}}^{{}^{\kappa-{\rm n.t.}}}$ exists 
{\rm (}in ${\mathbb{C}}^n${\rm )} at ${\mathscr{L}}^{n-1}$-a.e. point in ${\mathbb{R}}^{n-1}$.
\end{enumerate}

\noindent Then, with $e_n:=(0,\dots,0,1)\in{\mathbb{R}}^n$ and ``dot" denoting the
standard inner product in ${\mathbb{R}}^n$,
%
\begin{equation}\label{eqn:div-form}
{}_{(\mathscr{C}_b^\infty({\mathbb{R}}^n_{+}))^\ast}\big\langle{\rm div}\,\vec{F},1
\big\rangle_{\mathscr{C}_b^\infty({\mathbb{R}}^n_{+})}
=-\int_{{\mathbb{R}}^{n-1}}e_n\cdot
\big(\vec F\,\big|^{{}^{\kappa-{\rm n.t.}}}_{\partial{\mathbb{R}}^n_{+}}\bigr)\,d{\mathscr{L}}^{n-1}.
\end{equation}
\end{theorem}

\begin{lemma}\label{Lgav-TeD}
Assume that $L$ is an $M\times M$ constant complex coefficient system as in \eqref{L-def} and \eqref{Def-ELLa.INTR}, 
and fix an aperture parameter $\kappa>0$. Pick two points $x^\star_j\in{\mathbb{R}}^{n}_{+}$ 
and two radii $r_j\in(0,\infty)$, $j\in\{1,2\}$, satisfying 
%
\begin{equation}\label{CPTs}
r_j<{\rm dist}\,(x^\star_j,\partial{\mathbb{R}}^n_{+})\frac{\kappa}{\sqrt{1+\kappa^2}}
\,\,\text{ for }\,\,j\in\{1,2\},\,\,\text{ and }\,\,|x^\star_1-x^\star_2|>r_1+r_2, 
\end{equation}
%
then abbreviate $K_j:=\overline{B(x^\star_j,r_j)}$ for $j\in\{1,2\}$.

Next, consider two vector-valued functions
%
\begin{equation}\label{Tvav-TREINRv}
u=\big(u_{\alpha}\big)_{1\leq\alpha\leq M}\in\big[W^{1,1}_{\rm loc}({\mathbb{R}}^n_{+})\big]^M,\quad
v=\big(v_{\alpha}\big)_{1\leq\alpha\leq M}\in\big[W^{1,1}_{\rm loc}({\mathbb{R}}^n_{+})\big]^M
\end{equation}
%
satisfying 
%
\begin{align}\label{hBbac.1}
& {\mathcal{N}}_\kappa^{\,\mathbb{R}^{n}_{+}\setminus K_2}v,\,
{\mathcal{N}}_\kappa^{\,\mathbb{R}^{n}_{+}\setminus K_2}(\nabla v)\in L^1_{\rm loc}({\mathbb{R}}^{n-1}),\quad
L^{\top}v\in\big[{\mathcal{E}}'_{K_2}(\mathbb{R}^{n}_{+})\big]^M,
\\[4pt]
& v\big|^{{}^{\kappa-{\rm n.t.}}}_{\partial{\mathbb{R}}^n_{+}}=0
\,\,\text{ a.e. in }\,\,{\mathbb{R}}^{n-1},\quad
\big(\nabla v\big)\big|_{\partial\mathbb{R}^{n}_{+}}^{{}^{\kappa-{\rm n.t.}}}
\,\,\text{ exists a.e. in }\,\,{\mathbb{R}}^{n-1},
\label{hBbac.2}
\\[4pt]
\label{DiBBa.LLap}
& Lu\in\big[{\mathcal{E}}'_{K_1}(\mathbb{R}^{n}_{+})\big]^M,\quad
u\big|_{\partial\mathbb{R}^{n}_{+}}^{{}^{\kappa-{\rm n.t.}}}\,\,\text{ exists a.e. in }\,\,{\mathbb{R}}^{n-1},
\\[4pt]
& \text{and }\,\,
\big({\mathcal{N}}^{\,\mathbb{R}^{n}_{+}\setminus K_1}_\kappa u\big)\cdot
\big({\mathcal{N}}_\kappa^{\,\mathbb{R}^{n}_{+}\setminus K_2}(\nabla v)\big)\in L^1({\mathbb{R}}^{n-1}).
\label{DiBBa.LLap.BBBB}
\end{align}
%
Finally, select two scalar functions, 
%
\begin{equation}\label{Ua-eNBVVatGG}
\begin{array}{c}
\phi\in{\mathscr{C}}^\infty_c({\mathbb{R}}^n_{+})
\,\,\text{ such that $\phi\equiv 1$ near $K_1$ and $\phi\equiv 0$ near $K_2$},
\\[6pt]
\psi\in{\mathscr{C}}^\infty_c({\mathbb{R}}^n_{+})
\,\,\text{ such that $\psi\equiv 1$ near $K_2$ and $\psi\equiv 0$ near $K_1$}.
\end{array}
\end{equation}

Then
%
\begin{align}\label{Ua-eDPa4.L}
{}_{[{\mathcal{E}}'(\mathbb{R}^{n}_{+})]^M}\big\langle Lu,\phi v
\big\rangle_{[{\mathcal{E}}(\mathbb{R}^{n}_{+})]^M}
=&\,
{}_{[{\mathcal{E}}'(\mathbb{R}^{n}_{+})]^M}\big\langle L^{\top}v,
\psi u\big\rangle_{[{\mathcal{E}}(\mathbb{R}^{n}_{+})]^M}
\nonumber\\[4pt]
&\,+\int_{\mathbb{R}^{n-1}}
\Big(u_\alpha\big|_{\partial\mathbb{R}^{n}_{+}}^{{}^{\kappa-{\rm n.t.}}}\Big) 
a^{\beta\alpha}_{nn}\,\Big[\big(\partial_n v_\beta\big)\Big]
\Big|_{\partial\mathbb{R}^{n}_{+}}^{{}^{\kappa-{\rm n.t.}}}\,d{\mathscr{L}}^{n-1}.
\end{align}
\end{lemma}

\begin{proof}
Note that the ellipticity of $L^{\top}$ (entailed by that of $L$) and the 
fact that $v\in\big[L^1_{\rm loc}({\mathbb{R}}^n_{+})\big]^M$ satisfies
$L^{\top}v\in\big[{\mathcal{E}}'_{K_2}(\mathbb{R}^{n}_{+})\big]^M$ imply that 
%
\begin{equation}\label{Di-UPs-4}
v\in\big[{\mathscr{C}}^\infty(\mathbb{R}^{n}_{+}\setminus K_2)\big]^M\,\,\text{ and }\,\,
L^{\top}v=0\,\,\text{ pointwise in }\,\,\mathbb{R}^{n}_{+}\setminus K_2.
\end{equation}
%
Likewise, the conditions on $u$ give  
%
\begin{equation}\label{Di-UabTRF.a}
u\in\big[{\mathscr{C}}^\infty(\mathbb{R}^{n}_{+}\setminus K_1)\big]^M\,\,\text{ and }\,\,
Lu=0\,\,\text{ pointwise in }\,\,\mathbb{R}^{n}_{+}\setminus K_1.
\end{equation}
%
We shall first establish formula \eqref{Ua-eDPa4.L} under the additional assumption that 
%
\begin{equation}\label{Di-UabTRF.b}
\begin{array}{c}
\big({\mathcal{N}}_\kappa^{\,\mathbb{R}^{n}_{+}\setminus K_2}v\big)\cdot
\big({\mathcal{N}}^{\,\mathbb{R}^{n}_{+}\setminus K_1}_\kappa(\nabla u)\big)\in L^1({\mathbb{R}}^{n-1})
\,\,\text{ and }\,\,(\nabla u)\Big|_{\partial\mathbb{R}^{n}_{+}}^{{}^{\kappa-{\rm n.t.}}}\,\,
\text{ exists a.e. in }\,\,{\mathbb{R}}^{n-1}.
\end{array}
\end{equation}
%
Assuming that this is the case define 
%
\begin{equation}\label{Yab-HBBaR46}
\vec{F}:=\Big(v_{\alpha}\,a^{\alpha\beta}_{jk}\,\partial_k u_\beta
-u_\alpha a^{\beta\alpha}_{kj}\partial_k v_{\beta}\Big)_{1\le j\le n}
\,\,\,\text{ in }\,\,{\mathbb{R}}^n_{+}.
\end{equation}
%
Then \eqref{Yab-HBBaR46} entails 
%
\begin{equation}\label{Yab-HBAAA.h5}
\big|\vec{F}\big|\leq C\big(|v||\nabla u|+|u|\big|\nabla v\big|\big)
\,\,\text{ a.e. in }\,\,{\mathbb{R}}^n_{+}.
\end{equation}
%
Note that $|v|$ is locally absolutely integrable (by \eqref{Tvav-TREINRv}), 
while $|\nabla u|$ is bounded in a neighborhood of $K_2$ (by \eqref{Di-UabTRF.a} 
and the fact that $K_1\cap K_2=\varnothing$), hence the product $|v||\nabla u|$ is 
absolutely integrable in a neighborhood of $K_2$. Away from $K_2$, we have that 
$|v|$ is locally bounded (by \eqref{Di-UPs-4}), while $|\nabla u|$ is locally 
integrable (by \eqref{Tvav-TREINRv}). Thus, the product $|v||\nabla u|$ is also 
locally absolutely integrable away from $K_2$. In summary, 
$|v||\nabla u|\in L^1_{\rm loc}({\mathbb{R}}^n_{+})$. In a similar manner, we also 
obtain that $|u||\nabla v|\in L^1_{\rm loc}({\mathbb{R}}^n_{+})$, 
hence ultimately, 
%
\begin{equation}\label{Yab-HBAAA.aF}
\vec{F}\in\big[L^1_{\rm loc}({\mathbb{R}}^n_{+})\big]^n.
\end{equation}
%
In particular, it makes sense to consider ${\rm div}\vec{F}$ in the sense of 
distributions in ${\mathbb{R}}^n_{+}$. In fact, a direct calculation gives
%
\begin{equation}\label{div-F.Req}
{\rm div}\vec{F}=-u_{\alpha}\,(L^{\top}v)_\alpha+v_{\alpha}\,(Lu)_\alpha
\,\,\,\text{ in }\,\,{\mathcal{D}}'({\mathbb{R}}^n_{+}).
\end{equation}
%
Consequently, if $K_\star:=K_1\cup K_2$ then $K_\star$ is a compact 
subset of ${\mathbb{R}}^n_{+}$ and, as seen from \eqref{div-F.Req}, the last membership in \eqref{hBbac.1},
the membership in \eqref{DiBBa.LLap}, \eqref{Di-UPs-4}-\eqref{Di-UabTRF.a}, we have
%
\begin{equation}\label{div-F.Req.3}
{\rm div}\vec{F}\in{\mathcal{E}}'_{K_\star}({\mathbb{R}}^n_{+}). 
\end{equation}
%
Furthermore, \eqref{Yab-HBAAA.h5} gives
%
\begin{equation}\label{Yab-HBna-TAn}
{\mathcal{N}}^{\,\mathbb{R}^{n}_{+}\setminus K_\star}_{\kappa}\vec{F}
\leq C\Big({\mathcal{N}}^{\,\mathbb{R}^{n}_{+}\setminus K_2}_{\kappa}v\cdot
{\mathcal{N}}^{\,\mathbb{R}^{n}_{+}\setminus K_1}_{\kappa}(\nabla u) 
+{\mathcal{N}}^{\,\mathbb{R}^{n}_{+}\setminus K_1}_{\kappa}u
\cdot{\mathcal{N}}^{\,\mathbb{R}^{n}_{+}\setminus K_2}_{\kappa}(\nabla v)\Big)
\,\,\text{ in }\,\,{\mathbb{R}}^{n-1}.
\end{equation}
%
From \eqref{Yab-HBna-TAn}, the first working assumption made in \eqref{Di-UabTRF.b}, and \eqref{DiBBa.LLap.BBBB}
we therefore obtain (also bearing in mind that the nontangential maximal function is lower-semicontinuous, hence 
measurable)
%
\begin{equation}\label{Yab-HBna-TAn.2}
{\mathcal{N}}^{\,\mathbb{R}^{n}_{+}\setminus K_\star}_{\kappa}\vec{F}\in L^1({\mathbb{R}}^{n-1}).
\end{equation}
%
Finally, from \eqref{Yab-HBBaR46}, \eqref{hBbac.2}, the second working assumption made in \eqref{Di-UabTRF.b}, 
and the middle condition in \eqref{DiBBa.LLap}, we see that
$\vec{F}\,\big|^{{}^{\kappa-{\rm n.t.}}}_{\partial\mathbb{R}^n_+}$
exists a.e. in ${\mathbb{R}}^{n-1}$ and, in fact, 
%
\begin{equation}\label{KnaJHcwkfc}
\vec{F}\,\Big|^{{}^{\kappa-{\rm n.t.}}}_{\partial\mathbb{R}^n_+}
=-\Big(\Big(u_\alpha\Big|^{{}^{\kappa-{\rm n.t.}}}_{\partial\mathbb{R}^n_+}\Big)
a^{\beta\alpha}_{kj}\Big(\big(\partial_k v_{\beta}\big)
\Big|^{{}^{\kappa-{\rm n.t.}}}_{\partial\mathbb{R}^n_+}\Big)_{1\leq j\le n}
\,\,\text{ a.e. on }\,\,{\mathbb{R}}^{n-1}.
\end{equation}
%
Then Theorem~\ref{theor:div-thm} applies and on account of \eqref{div-F.Req}
and \eqref{KnaJHcwkfc} yields, for any two given functions $\psi,\phi$ 
as in \eqref{Ua-eNBVVatGG},
%
\begin{align}\label{Uahg-ihGPPa4}
&\hskip -0.40in
-{}_{[{\mathcal{E}}'(\mathbb{R}^{n}_{+})]^M}\big\langle L^{\top}v,
\psi u\big\rangle_{[{\mathcal{E}}(\mathbb{R}^{n}_{+})]^M}
+{}_{[{\mathcal{E}}'(\mathbb{R}^{n}_{+})]^M}\big\langle Lu\,,\,
\phi\,v\big\rangle_{[{\mathcal{E}}(\mathbb{R}^{n}_{+})]^M}
\nonumber\\[4pt]
&\hskip 0.80in
={}_{(\mathscr{C}_b^\infty(\mathbb{R}^n_+))^*}\big\langle{\rm div}\vec{F},1
\big\rangle_{\mathscr{C}_b^\infty(\mathbb{R}^n_+)} 
=-\int_{\mathbb{R}^{n-1}}e_n\cdot
\big(\vec F\,\big|^{{}^{\kappa-{\rm n.t.}}}_{\partial\mathbb{R}^n_+}\bigr)\,d{\mathscr{L}}^{n-1}
\nonumber\\[4pt]
&\hskip 0.80in
=\int_{\mathbb{R}^{n-1}}
\Big(u_\alpha\big|_{\partial\mathbb{R}^{n}_{+}}^{{}^{\kappa-{\rm n.t.}}}\Big)a^{\beta\alpha}_{kn}
\Big[\big(\partial_k v_{\beta}\big)
\Big]\Big|_{\partial\mathbb{R}^{n}_{+}}^{{}^{\kappa-{\rm n.t.}}}\,d{\mathscr{L}}^{n-1}.
\end{align}

At this stage, we claim that 
%
\begin{equation}\label{Ua-nabUGC4D}
\begin{array}{c}
\Big[\big(\partial_k v_\beta\big)\Big]
\Big|_{\partial\mathbb{R}^{n}_{+}}^{{}^{\kappa-{\rm n.t.}}}=0
\quad\text{ a.e. on }\,\,\partial{\mathbb{R}}^n_{+}\equiv\mathbb{R}^{n-1},
\\[10pt]
\forall\,k\in\{1,\dots,n-1\}\,\,\text{ and }\,\,\forall\,\beta\in\{1,\dots,M\}.
\end{array}
\end{equation}
%
To see that this is the case, fix $k\in\{1,\dots,n-1\}$ along with 
$\beta\in\{1,\dots,M\}$, and pick an arbitrary function 
%
\begin{equation}\label{Ua-nabUGC4D.2}
\varphi\in{\mathscr{C}}^\infty_c({\mathbb{R}}^n)\,\,\text{ such that }\,\,
K_2\cap{\rm supp}\,\varphi=\emptyset. 
\end{equation}
%
Next, consider the vector field
%
\begin{equation}\label{Ua-nabUGC4D.3}
\vec{H}:=\Big(\big(\partial_k v_\beta\big)\varphi+v_\beta\partial_k\varphi\Big)e_n
-\Big(\big(\partial_n v_{\beta}\big)\varphi+v_{\beta}\partial_n\varphi\Big)e_k
\,\,\,\text{ in }\,\,{\mathbb{R}}^n_{+}.
\end{equation}
%
Note that 
%
\begin{equation}\label{Ua-nabUGC4D.4}
\vec{H}\in\big[{\mathscr{C}}^\infty({\mathbb{R}}^n_{+})\big]^n\,\,\text{ and }\,\,
{\rm div}\vec{H}=0\,\,\text{ in }\,\,{\mathcal{D}}'({\mathbb{R}}^n_{+}),
\end{equation}
%
by \eqref{Ua-nabUGC4D.3}, \eqref{Ua-nabUGC4D.2}, and \eqref{Di-UPs-4}. 
Also, since ${\mathcal{N}}_{\kappa}\varphi$, ${\mathcal{N}}_{\kappa}(\nabla\varphi)$ 
are bounded functions with compact support in ${\mathbb{R}}^{n-1}$, we have
%
\begin{equation}\label{Ua-nabUGC4D.5}
{\mathcal{N}}_{\kappa}\vec{H}\in L^1({\mathbb{R}}^{n-1})
\end{equation}
%
by \eqref{Ua-nabUGC4D.3}, \eqref{Ua-nabUGC4D.2}, and \eqref{hBbac.1}.
In addition, thanks to \eqref{hBbac.2}, 
%
\begin{equation}\label{Ua-nabUGC4D.6}
\vec{H}\Big|_{\partial\mathbb{R}^{n}_{+}}^{{}^{\kappa-{\rm n.t.}}}
=\big(\varphi\big|_{\partial\mathbb{R}^{n}_{+}}\big)
\Big[\big(\partial_k v_{\beta}\big)\Big]
\Big|_{\partial\mathbb{R}^{n}_{+}}^{{}^{\kappa-{\rm n.t.}}}\,e_n
-\big(\varphi\big|_{\partial\mathbb{R}^{n}_{+}}\big)
\Big[\big(\partial_n v_{\beta}\big)\Big]
\Big|_{\partial\mathbb{R}^{n}_{+}}^{{}^{\kappa-{\rm n.t.}}}\,e_k,
\end{equation}
%
at a.e. point on ${\mathbb{R}}^{n-1}$.
Granted this and keeping in mind that $k\in\{1,\dots,n-1\}$, 
we therefore obtain 
%
\begin{equation}\label{Ua-nabUGC4D.7}
e_n\cdot\Big(\vec{H}\Big|_{\partial\mathbb{R}^{n}_{+}}^{{}^{\kappa-{\rm n.t.}}}\Big)
=\big(\varphi\big|_{\partial\mathbb{R}^{n}_{+}}\big)
\Big[\big(\partial_k v_{\beta}\big)\Big]
\Big|_{\partial\mathbb{R}^{n}_{+}}^{{}^{\kappa-{\rm n.t.}}}
\,\,\text{ a.e. in }\,\,{\mathbb{R}}^{n-1}.
\end{equation}
%
Collectively, \eqref{Ua-nabUGC4D.4}-\eqref{Ua-nabUGC4D.6} ensure that 
Theorem~\ref{theor:div-thm} is applicable to the vector field $\vec{H}$. 
Based on the second condition in \eqref{Ua-nabUGC4D.4} and \eqref{Ua-nabUGC4D.7} 
we may therefore write 
%
\begin{align}\label{Phgav-Y54V}
0 & ={}_{(\mathscr{C}_b^\infty(\mathbb{R}^n_+))^*}\big\langle{\rm div}\vec{H},1
\big\rangle_{\mathscr{C}_b^\infty(\mathbb{R}^n_+)} 
=-\int_{\mathbb{R}^{n-1}}e_n\cdot
\big(\vec{H}\,\big|^{{}^{\kappa-{\rm n.t.}}}_{\partial\mathbb{R}^n_+}\big)\,d{\mathscr{L}}^{n-1}
\nonumber\\[4pt]
&=-\int_{\mathbb{R}^{n-1}}\big(\varphi\big|_{\partial\mathbb{R}^{n}_{+}}\big)
\Big[\big(\partial_k v_{\beta}\big)
\Big]\Big|_{\partial\mathbb{R}^{n}_{+}}^{{}^{\kappa-{\rm n.t.}}}\,d{\mathscr{L}}^{n-1}.
\end{align}
%
On the other hand, 
%
\begin{equation}\label{uag-tREW.1}
\Big\{\varphi\big|_{\partial\mathbb{R}^{n}_{+}}:\,
\varphi\in{\mathscr{C}}^\infty_c({\mathbb{R}}^n)\,\,
\text{ such that }\,\,K_2\cap{\rm supp}\,\varphi=\varnothing\Big\}
={\mathscr{C}}^\infty_c({\mathbb{R}}^{n-1})
\end{equation}
%
while from \eqref{hBbac.1}-\eqref{hBbac.2} we have 
%
\begin{equation}\label{ubav-KBav4fE}
\Big[\big(\partial_k v_{\beta}\big)
\Big]\Big|_{\partial\mathbb{R}^{n}_{+}}^{{}^{\kappa-{\rm n.t.}}}\in L^1_{\rm loc}({\mathbb{R}}^{n-1}).
\end{equation}
%
Together, \eqref{ubav-KBav4fE}, \eqref{Phgav-Y54V} and \eqref{uag-tREW.1} 
ultimately prove that $\big[\big(\partial_k v_{\beta}\big)\big]
\big|_{\partial\mathbb{R}^{n}_{+}}^{{}^{\kappa-{\rm n.t.}}}=0$ a.e. in ${\mathbb{R}}^{n-1}$, 
finishing the justification of the claim made in \eqref{Ua-nabUGC4D}. 
In turn, \eqref{Uahg-ihGPPa4} and \eqref{Ua-nabUGC4D} prove \eqref{Ua-eDPa4.L} 
under the additional assumption \eqref{Di-UabTRF.b}. 

To dispense with \eqref{Di-UabTRF.b}, fix $\phi,\psi$ as in \eqref{Ua-eNBVVatGG}.
For each $\varepsilon>0$ define $u^\varepsilon=(u^\varepsilon_\alpha)_{1\leq\alpha\leq M}$ by 
%
\begin{equation}\label{Nvav-T54.a}
u^\varepsilon(x):=u(x+\varepsilon e_n),\qquad\forall\,x\in{\mathbb{R}}^n_{+}.
\end{equation}
%
Then $u^\varepsilon\in\big[W^{1,1}_{\rm loc}({\mathbb{R}}^n_{+})\big]^{M}$ and $u^\varepsilon$ extends to a  
${\mathscr{C}}^\infty$ function in a two-sided neighborhood of $\partial{\mathbb{R}}^n_{+}$ provided 
$\varepsilon>0$ is small enough. In particular, both 
$u^\varepsilon\big|_{\partial\mathbb{R}^{n}_{+}}^{{}^{\kappa-{\rm n.t.}}}$ and 
$(\nabla u^\varepsilon)\big|_{\partial\mathbb{R}^{n}_{+}}^{{}^{\kappa-{\rm n.t.}}}$ exist a.e. in ${\mathbb{R}}^{n-1}$.
Also, if $0<\varepsilon<{\rm dist}\,(K_1,\partial{\mathbb{R}}^n_{+})$ it follows that 
$K_1-\varepsilon e_n$ is a compact subset of ${\mathbb{R}}^n_{+}$ and 
$Lu^\varepsilon\in\big[{\mathcal{E}}'_{K_1-\varepsilon e_n}(\mathbb{R}^{n}_{+})\big]^M$.

Moving on, fix $\kappa_o\in(0,\kappa)$ such that 
%
\begin{equation}\label{CPTs.i}
r_j<{\rm dist}\,(x^\star_j,\partial{\mathbb{R}}^n_{+})\frac{\kappa_o}{\sqrt{1+\kappa^2_o}}
\,\,\text{ for }\,\,j\in\{1,2\}.
\end{equation}
%
Let us now choose $x'\in{\mathbb{R}}^{n-1}\equiv\partial{\mathbb{R}}^n_{+}$ 
with the property that $\big(v\big|^{{}^{\kappa-{\rm n.t.}}}_{\partial{\mathbb{R}}^n_{+}}\big)(x')=0$.
We claim that 
%
\begin{equation}\label{Cr654.A-HM.1}
|v(y)|\leq Cy_n\cdot\big({\mathcal{N}}^{\,\mathbb{R}^{n}_{+}\setminus K_2}_{\kappa}(\nabla v)\big)(x')
\,\,\text{ for each }\,\,y=(y',y_n)\in\Gamma_{\kappa_o}(x')\setminus K_2.
\end{equation}
% 
To see this, let $y$ be as in \eqref{Cr654.A-HM.1} and choose a rectifiable path 
$\gamma:[0,1]\to\overline{{\mathbb{R}}^n_{+}}$ joining $(x',0)$ with $y$, whose length is $\leq Cy_n$, 
and such that $\gamma((0,1))\subseteq\Gamma_{\kappa_o}(x')\setminus K_2$.
Then, for some constant $C\in(0,\infty)$ independent of $x'$ and $y$, we may estimate 
%
\begin{align}\label{Cr654.A-HM.1.xxx}
|v(y)| &=\Big|v(y)-\big(v\big|^{{}^{\kappa-{\rm n.t.}}}_{\partial{\mathbb{R}}^n_{+}}\big)(x')\Big|
=\Big|\int_0^1\frac{d}{dt}[v(\gamma(t))]\,dt\Big|
\nonumber\\[4pt]
&=\Big|\int_0^1(\nabla v)(\gamma(t))\cdot\gamma'(t)\,dt\Big|
\leq\Big(\sup_{\xi\in\gamma((0,1))}|(\nabla v)(\xi)|\Big)\int_0^1|\gamma'(t)|\,dt
\nonumber\\[4pt]
&\leq Cy_n\cdot\big({\mathcal{N}}^{\,\mathbb{R}^{n}_{+}\setminus K_2}_{\kappa_o}(\nabla v)\big)(x')
\leq Cy_n\cdot\big({\mathcal{N}}^{\,\mathbb{R}^{n}_{+}\setminus K_2}_{\kappa}(\nabla v)\big)(x'),
\end{align}
%
using the Fundamental Theorem of Calculus, Chain Rule, and \eqref{NT-Fct.23}. This establishes \eqref{Cr654.A-HM.1}.

For the remainder of the proof assume that the parameter $\varepsilon>0$ is sufficiently small, say 
$\varepsilon<\frac{1}{2}\,{\rm dist}\,(K_1,\partial{\mathbb{R}}^n_{+})$, and define 
%
\begin{align}\label{Cr654.A-HM.2-FDS.aTgab}
K_1^\varepsilon:=\{z\in{\mathbb{R}}^n:{\rm dist}\,(z,K_1-\varepsilon e_n)\leq\varepsilon\}.
\end{align}
%
This is a compact subset of ${\mathbb{R}}^n_{+}$. Next, choose some sufficiently small 
parameter $a\in(0,1)$. Specifically, start with some parameter $a$ satisfying 
%
\begin{align}\label{Cr654.A-HM.2-FDS.aaa}
0<a<\frac{\kappa-\kappa_o}{1+\kappa}
\end{align}
%
and which is small enough so that 
%
\begin{align}\label{Cr654.A-HM.2-FDS.a344}
B(y,a\cdot y_n)\cap(K_1-\varepsilon e_n)=\varnothing\,\,\text{ for each }\,\,
y\in{\mathbb{R}}^n_{+}\setminus K_1^\varepsilon.
\end{align}
%
To see why this may be arranged, pick $R:=2\,\sup\{z_n:\,z=(z',z_n)\in K_1-\varepsilon e_n\}$
then take $a<\min\big\{1/2\,,\,\varepsilon/R\big\}$. Pick an arbitrary point 
$y=(y',y_n)\in{\mathbb{R}}^n_{+}\setminus K_1^\varepsilon$. On the one hand, if $y_n\geq R$ then 
$B(y,a\cdot y_n)$ lies above the strip $\{z=(z',z_n)\in{\mathbb{R}}^n_{+}:z_n<(1-a)R\}$ 
and the latter strip contains the set $K_1-\varepsilon e_n$. Thus, the disjointness condition in 
\eqref{Cr654.A-HM.2-FDS.a344} holds in this case. On the other hand, if $y_n<R$ then 
$a\cdot y_n<a\cdot R<\varepsilon$, and since $B(y,\varepsilon)$ does not intersect 
$K_1-\varepsilon e_n$ (given that ${\rm dist}\,(y,K_1-\varepsilon e_n)>\varepsilon$
since $y\in{\mathbb{R}}^n_{+}\setminus K_1^\varepsilon$), the disjointness condition
in \eqref{Cr654.A-HM.2-FDS.a344} holds in this case as well. 

Having chosen $a$ as in \eqref{Cr654.A-HM.2-FDS.aaa} so that \eqref{Cr654.A-HM.2-FDS.a344} holds, we now fix an arbitrary point 
$y\in\Gamma_{\kappa_o}(x')\setminus K_1^\varepsilon$. Observe that having $z=(z',z_n)\in B(y,a\cdot y_n)$ entails 
%
\begin{align}\label{Cr654.A-HM.2-FDS.1}
y_n\leq z_n+|z-y|<z_n+a\cdot y_n\Longrightarrow y_n<(1-a)^{-1}z_n,
\end{align}
%
which, bearing in mind the current location of $y$, permits us to conclude that  
%
\begin{align}\label{Cr654.A-HM.2-FDS.2}
|z'-x'| &\leq|z'-y'|+|y'-x'|\leq|z-y|+\kappa_o\cdot y_n<a\cdot y_n+\kappa_o\cdot y_n
\nonumber\\[6pt]
&=(\kappa_o+a)y_n<\frac{\kappa_o+a}{1-a}z_n<\kappa z_n,
\end{align}
%
where the last inequality is guaranteed by \eqref{Cr654.A-HM.2-FDS.aaa}.
From \eqref{Cr654.A-HM.2-FDS.2} we see that $z\in\Gamma_{\kappa}(x')$.
Together with \eqref{Cr654.A-HM.2-FDS.a344}, this proves that 
%
\begin{align}\label{Cr654.A-HM.2-FDS.a344.ag}
B(y,a\cdot y_n)\subseteq\Gamma_{\kappa}(x')\setminus(K_1-\varepsilon e_n)\,\,\text{ for each }\,\,
y\in\Gamma_{\kappa_o}(x')\setminus K_1^\varepsilon.
\end{align}
%
Using interior estimates (cf. Theorem~\ref{ker-sbav}) for the function $u^\varepsilon$ we may then write 
%
\begin{align}\label{Cr654.A-HM.2-bis}
|(\nabla u^\varepsilon)(y)| &\leq\frac{C}{y_n}\meanint_{B(y,a\cdot y_n)}|u^\varepsilon(z)|\,dz
\nonumber\\[6pt]
&\leq Cy_n^{-1}\cdot
\sup_{z\in\Gamma_{\kappa}(x')\setminus(K_1-\varepsilon e_n)}|u^\varepsilon(z)|
\nonumber\\[6pt]
&\leq Cy_n^{-1}\cdot\big({\mathcal{N}}_{\kappa}^{\,\mathbb{R}^{n}_{+}\setminus (K_1-\varepsilon e_n)}u^\varepsilon\big)(x')
\nonumber\\[6pt]
&\leq Cy_n^{-1}\cdot\big({\mathcal{N}}_{\kappa}^{\,\mathbb{R}^{n}_{+}\setminus K_1}u\big)(x'),
\end{align}
%
where $C=C(L,n,a)\in(0,\infty)$. Combining \eqref{Cr654.A-HM.1} with \eqref{Cr654.A-HM.2-bis} then gives
%
\begin{align}\label{Cr654.A-HM.2-b22}
\big({\mathcal{N}}^{\,\mathbb{R}^{n}_{+}\setminus (K_1^\varepsilon\cup K_2)}_{\kappa_o}
(|v||\nabla u^\varepsilon|)\big)(x')
\leq C\big({\mathcal{N}}^{\,\mathbb{R}^{n}_{+}\setminus K_2}_{\kappa}(\nabla v)\big)(x')
\big({\mathcal{N}}^{\,\mathbb{R}^{n}_{+}\setminus K_1}_{\kappa}u\big)(x').
\end{align}
%
Also, it is clear that 
%
\begin{align}\label{Cr654.A-HM.2-b22.bb}
\big({\mathcal{N}}^{\,\mathbb{R}^{n}_{+}\setminus (K_1^\varepsilon\cup K_2)}_{\kappa}(|\nabla v||u^\varepsilon|)\big)(x')
\leq C\big({\mathcal{N}}^{\,\mathbb{R}^{n}_{+}\setminus K_2}_{\kappa}(\nabla v)\big)(x')
\big({\mathcal{N}}^{\,\mathbb{R}^{n}_{+}\setminus K_1}_{\kappa}u\big)(x').
\end{align}
%
In concert with \eqref{DiBBa.LLap.BBBB}, these ultimately show that if we now define 
%
\begin{equation}\label{Yab-HBBaR46-bis}
\vec{F}^\varepsilon:=\Big(v_{\alpha}\,a^{\alpha\beta}_{jk}\,\partial_k u^\varepsilon_\beta
-u^\varepsilon_\alpha a^{\beta\alpha}_{kj}\partial_k v_{\beta}\Big)_{1\le j\le n}
\,\,\,\text{ in }\,\,{\mathbb{R}}^n_{+},
\end{equation}
%
then 
%
\begin{equation}\label{Yab-HBna-TAn.2-bis}
{\mathcal{N}}^{\,\mathbb{R}^{n}_{+}\setminus(K_1^\varepsilon\cup K_2)}_{\kappa_o}\vec{F}^\varepsilon\in L^1({\mathbb{R}}^{n-1}).
\end{equation}
%
Granted this, the same type of argument as in the first part of the proof 
relying on Theorem~\ref{theor:div-thm} (now applied to $\vec{F}^\varepsilon$ in place of $\vec{F}$, with 
$K_1^\varepsilon\cup K_2$ in place of $K_\star$, and with for the aperture parameter $\kappa_o$) gives that
%
\begin{align}\label{Ua-eDPa4.LAC}
{}_{[{\mathcal{E}}'(\mathbb{R}^{n}_{+})]^M}\big\langle Lu^\varepsilon,\phi v
\big\rangle_{[{\mathcal{E}}(\mathbb{R}^{n}_{+})]^M}
=&\,{}_{[{\mathcal{E}}'(\mathbb{R}^{n}_{+})]^M}\big\langle L^{\top}v,
\psi u^\varepsilon\big\rangle_{[{\mathcal{E}}(\mathbb{R}^{n}_{+})]^M}
\nonumber\\[4pt]
&\,+\int_{\mathbb{R}^{n-1}}
\Big(u^\varepsilon_\alpha\big|_{\partial\mathbb{R}^{n}_{+}}^{{}^{\kappa-{\rm n.t.}}}\Big) 
a^{\beta\alpha}_{nn}\,\Big[\big(\partial_n v_\beta\big)\Big]
\Big|_{\partial\mathbb{R}^{n}_{+}}^{{}^{\kappa-{\rm n.t.}}}\,d{\mathscr{L}}^{n-1}.
\end{align}

There remains to eliminate $\varepsilon$. To this end, observe that 
%
\begin{equation}\label{Ua-eDPa4.LAC.2}
{}_{[{\mathcal{E}}'(\mathbb{R}^{n}_{+})]^M}\big\langle Lu^\varepsilon,\phi v
\big\rangle_{[{\mathcal{E}}(\mathbb{R}^{n}_{+})]^M}
={}_{[{\mathcal{E}}'(\mathbb{R}^{n}_{+})]^M}\big\langle 
Lu\,,\,\phi(\cdot-\varepsilon e_n)v(\cdot-\varepsilon e_n)
\big\rangle_{[{\mathcal{E}}(\mathbb{R}^{n}_{+})]^M},
\end{equation}
%
and that (as seen from the smoothness properties of the functions involved)
%
\begin{equation}\label{naIaPa.2}
\lim\limits_{\varepsilon\to 0^{+}}
\big[\phi(\cdot-\varepsilon e_n)v(\cdot-\varepsilon e_n)\big]=\phi\,v
\,\,\text{ and }\,\,\lim\limits_{\varepsilon\to 0^{+}}\big[\psi u^\varepsilon\big]
=\psi\,u\,\,\text{ in }\,\,{\mathcal{E}}(\mathbb{R}^{n}_{+}).
\end{equation}
%
Hence, 
%
\begin{equation}\label{naIaPa.2.iiia}
\lim\limits_{\varepsilon\to 0^{+}}
{}_{[{\mathcal{E}}'(\mathbb{R}^{n}_{+})]^M}\big\langle Lu^\varepsilon,\phi v
\big\rangle_{[{\mathcal{E}}(\mathbb{R}^{n}_{+})]^M}
={}_{[{\mathcal{E}}'(\mathbb{R}^{n}_{+})]^M}\big\langle Lu,\phi v
\big\rangle_{[{\mathcal{E}}(\mathbb{R}^{n}_{+})]^M}
\end{equation}
%
and
%
\begin{equation}\label{naIaPa.2.iiib}
\lim\limits_{\varepsilon\to 0^{+}}
{}_{[{\mathcal{E}}'(\mathbb{R}^{n}_{+})]^M}\big\langle L^{\top}v,
\psi u^\varepsilon\big\rangle_{[{\mathcal{E}}(\mathbb{R}^{n}_{+})]^M}
={}_{[{\mathcal{E}}'(\mathbb{R}^{n}_{+})]^M}\big\langle L^{\top}v,
\psi u\big\rangle_{[{\mathcal{E}}(\mathbb{R}^{n}_{+})]^M}.
\end{equation}
%
Moreover, we have $|u(\cdot+\varepsilon e_n)|\leq{\mathcal{N}}^{\,\mathbb{R}^{n}_{+}\setminus K_1}_\kappa u$ 
on $\mathbb{R}^{n-1}$, and for each $\alpha\in\{1,\dots,M\}$ the last condition in \eqref{DiBBa.LLap} implies that 
%
\begin{equation}\label{naIaPa.3}
\lim\limits_{\varepsilon\to 0^{+}}u_\alpha(\cdot+\varepsilon e_n)
=u_\alpha\big|_{\partial\mathbb{R}^{n}_{+}}^{{}^{\kappa-{\rm n.t.}}}
\,\,\text{ a.e. in }\,\,\mathbb{R}^{n-1}.
\end{equation}
%
Consequently, formula \eqref{Ua-eDPa4.L} follows by letting $\varepsilon\to 0^{+}$
in \eqref{Ua-eDPa4.LAC}, on account of \eqref{Ua-eDPa4.LAC.2}-\eqref{naIaPa.3}, 
\eqref{DiBBa.LLap.BBBB}, and Lebesgue's Dominated Convergence Theorem.
\end{proof}

\end{document}
